\subsection{Methods for Computing Economic Equilibrium}
 

I will attempt to contextualize DEQNs \parencite{Azinovic2022DEQN} in the literature on economic equilibrium computation methods using Machine Learning (ML), emphasizing its relationship with the broader literature on solution algorithms that solve the kinds of problem formulations common in economics. I have found that this connection is often overlooked in literature reviews where DEQNs have been utilized, especially concerning its connection with Reinforcement Learning (RL). For a more comprehensive understanding of computational methods in economics, I suggest referring to the detailed survey on numerical methods in economics by \textcite{Maliar2014numerical_methods} and the survey on solution methods for DSGE models in \textcite{Villaverde2016survey_solution_dsge}.

Smolyak sparse grids, a widely used method originally introduced by the Russian mathematician \textcite{smolyak1963quadrature} have been frequently applied in an economic context, where several papers have demonstrated their efficacy in managing high-dimensional state spaces. For instance, \textcite{KruegerKubler2004} employed Smolyak sparse grids to tackle the computational challenges in solving overlapping generations (OLG) models with stochastic production. \textcite{WinschelKratzig2010} applied Smolyak sparse grids in a Bayesian estimation framework for nonlinear dynamic models. 

\textcite{BrummScheidegger2017sparse} introduced adaptive sparse grids. This advancement dynamically adjusts grid density in response to model complexity, allowing for even more precise and computationally efficient solutions to high-dimensional dynamic economic models. However, in their experiments, only going from solving a model with one asset to two assets was prohibitively difficult for them. The problem is that control variables need to be approximated on a hypercubic domain\footnote{Here, referring to the multi-dimensional grid representing the possible states of the system. Each axis represents a state variable, and the grid points represent the combinations of state variable values for which control needs to be computed.}. 

Traditional methods can struggle with dynamic models with stochasticity, high-dimensional state spaces, nonlinearities, and irregular state space geometries \parencite{Azinovic2022DEQN}. Grid-free simulation-based methods can be used to overcome these issues. A paper by \textcite{JuddMailar2011} highlights the numerical stability and accuracy of stochastic simulation approaches for solving dynamic models, indicating that simulation-based approaches outperform grid-based methods in certain contexts. \textcite{MaliarMailar2015} merge simulation with projection methods to tackle high-dimensional problems, showcasing the flexibility and potential of combining these techniques to overcome the curse of dimensionality. \textcite{ScheideggerBilionis2019} use a ML framework that leverages Gaussian processes and active subspaces to compute global solutions to dynamic stochastic economic models on irregular state spaces. This framework notably improves the ability to resolve local features and perform uncertainty quantification in high-dimensional settings, offering a significant advancement over traditional grid-based approaches.

Deep Neural Network (DNN) applications are becoming increasingly popular in algorithms used to estimate equilibrium in large economic models\footnote{See \textcite{Goodfellow2016} for an introduction to deep learning.}. While tangential to the main focus of this thesis, it is also noteworthy that DNNs have demonstrated capability in overcoming the curse of dimensionality in approximating partial differential equations (PDEs). For instance, \textcite{Jentzen2019diffs, Grohs2023diffs}, approximate Kolmogorov PDEs and Black-Sholes PDEs, respectively. 

Additionally, a paper by \textcite{becker2020deep} demonstrates that DNNs can approximate optimal stopping times to any given degree of accuracy, drawing inspiration from deep RL. Similarly, \textcite{gu2024pontryagin} explores the integration of DNNs with Pontryagin's Maximum Principle (PMP), effectively recovering optimality conditions and optimal action sequences. The approach was successful in various RL environments.

RL is a branch of ML that studies intertemporal optimization and, therefore, closely aligns with the problem statement in economic modeling and naturally has many parallels to the algorithms used for solving such problems. RL can be categorized into two main types: model-free and model-based. In model-free RL, the agent learns to make decisions based solely on observed states and rewards without any underlying model of the environment's dynamics. Conversely, in model-based RL, the agent makes decisions informed by an understanding of the environment's dynamics. Most economic algorithms fall under the latter, including the approach in \textcite{Azinovic2022DEQN} used in this thesis. For a thorough introduction to RL, I suggest referring to \textcite{Sutton2018intro_to_rl}, \textcite{Shakya2023} for recent advancements, and to \textcite{SpinningUp2018} for codes and references to state-of-the-art RL papers and codes. 

Close to the method in this thesis, \textcite{Duarte2018}, explicitly inspired by the RL literature, developed an algorithm for continuous-time finance that used DNNs combined with automatic differentiation and Ito's lemma to solve high-dimensional, nonlinear, continuous-time finance models. Another example of using DNNs in a continuous-time setting is presented in \textcite{Villaverde2023fin_frictions}. They use DNNs to solve for the perceived law of motion (PLM) in a model with heterogeneous agents.

The use of objective functions based on Euler equations, first-order conditions, and Bellman residuals is common in the literature. Examples of such approaches can be found in studies by \textcite{Valaitis2024, Maliar2021, Azinovic2022DEQN}, of which the last contains the solution algorithm used in the present thesis. All three studies present algorithms based on similar objectives and have similar algorithmic structures. Despite the similarities, \textcite{Valaitis2024} claims to do supervised learning, \textcite{Azinovic2022DEQN} claims it uses unsupervised learning, and  \textcite{Maliar2021} claims to be related to supervised, unsupervised and RL. 

Despite the varied descriptions provided by the authors, the core methodologies in these studies largely align with the learning paradigm of RL, or more specifically, model-based RL. This alignment is evident as all three approaches rely on iterative, knowledge-based policy updates driven by first-order conditions, Euler equations, or Bellman residuals. These elements are fundamental to RL, where interactions with an environment generate the data and learning signals, allowing us to approximate an optimal policy parameterized by a neural network. For more information on how RL distinguishes itself from supervised- and unsupervised learning, please refer to \textcite[p.~2]{Sutton2018intro_to_rl}. % I make this point as not everyone is aware they actually do RL. making this connection is crucial because of the advancements made in that literature

To my knowledge, this kind of learning algorithm has few applications to IAMs. As mentioned, two papers to date have used DEQNs as presented in \textcite{Azinovic2022DEQN}. \textcite{Folini2024ClimateInCE} used DEQNs to systematically recalibrate the climate emulator of DICE, ensuring their alignment with state-of-the-art climate data from CMIP5 \parencite{Taylor2012cmip5}. The paper by \textcite{friedl2023_deepUQ} employs DEQNs and Bayesian Learning to quantify uncertainty over the equilibrium climate sensitivity. Their method is applied to a high-dimensional IAM with stochastic climate tipping.

It is important to note that there have also been attempts to use model-free RL algorithms for IAMs. In one study, \textcite{strnard2019deepRLworld-earth} applied different variations of deep Q-learning \parencite{MnihVolodymyr2013PAwD} to the AYS model (as described in \textcite{Heitzig2018AYS}) and the model presented in \textcite{Nitzbon2017world-earth-mod}. Recently, \textcite{Wolf2023rl-to-iam}  explored more advanced RL algorithms, reward functions, and experimental setups to solve the AYS model.

This thesis makes several contributions to the literature. Firstly, it applies DEQNs to ACE, expanding the currently sparse literature that applies these methods to IAMs. Secondly, I provide a detailed explanation of the methods used, clarifying and strengthening the connection between DEQNs and RL. This can help future research in this area to draw on algorithms from the RL literature. 

Third, to the best of my knowledge, this work is one of the first attempts to validate the quality of DEQN solutions against analytical solutions. The analysis of the numerical results reveals areas where solution quality is suboptimal; the thesis specifically identifies the need for improved normalization processes in both pre- and post-processing stages as critical for enhancing model performance. 

Finally, the thesis is accompanied by a GitHub repository containing the codes and documentation for the results in the thesis. This repository is one of the few comprehensive, accessible resources available, which could significantly reduce the time and effort required for future research implementations.